% =====================================================================
%  1.2 Kvadratsetningene og faktorisering — ekstra utfordring
% =====================================================================

\begin{infoboks}{Velg metode selv}
Arket bygger på 1.2 Kvadratsetningene og faktorisering. Vis regning og begrunnelse.
Kontroller svaret ved å gange ut igjen eller sette inn et tall.
\end{infoboks}

\nivaahode{2}

\oppg{U1}{2}{\uh}
Gang ut og forenkle.
\[ (x+y+1)^2 \]

\oppg{U2}{2}{\uh}
Faktoriser fullstendig.
\begin{tasks}(2)
  \task $3x^3 - 27x$
  \task $12a^2 - 3b^2$
\end{tasks}

\nivaahode{3}

\oppg{U3}{3}{\uh}
Vis at $(n+1)^2 - n^2$ alltid er et oddetall, for alle heltall $n$.

\oppg{U4}{3}{\uh}
Finn $k$ slik at $x^2 + kx + 25$ er et fullstendig kvadrat.
Finn så $k$ slik at $x^2 + 6x + k$ er det.
Forklar hvorfor den første har to svar og den andre bare ett.

\oppg{U5}{3}{\prog}
Se på programmet.
\begin{kode}
for n in range(1, 11):
    v = (n + 1)**2 - n**2
    h = 2*n + 1
    print(n, v == h)
\end{kode}
\begin{tasks}(2)
  \task Hva skriver programmet ut? Kjør det, eller forklar uten å kjøre.
  \task Er dette et bevis for at $(n+1)^2 - n^2 = 2n + 1$ for alle $n$? Forklar.
\end{tasks}

\oppg{U6}{3}{\uh}
\begin{tasks}(2)
  \task Vis at $(m-d)(m+d) = m^2 - d^2$.
  \task Bruk a) til å regne ut $47 \cdot 53$ uten kalkulator.
  \task! Hvilke to hele tall med sum $100$ har størst produkt? Begrunn med $(m-d)(m+d) = m^2 - d^2$.
\end{tasks}

\oppg{U7}{3}{\uh}
\begin{tasks}(2)
  \task Faktoriser $x^4 - 16$ fullstendig.
  \task Forklar hvorfor $x^2 + 4$ ikke kan faktoriseres videre med reelle tall.
\end{tasks}

\oppg{U8}{3}{\uh}
Undersøk hvilke tall fra $1$ til $20$ som kan skrives som en differanse mellom to kvadrattall $0, 1, 4, 9, \ldots$
Hva har resten felles? Forklar med konjugatsetningen.

\begin{oppsummering}
Trekk ut felles faktor først. Velg deretter kvadratsetning eller konjugatsetning.
Gang ut igjen for å kontrollere svaret.
\end{oppsummering}
